\documentclass[11pt]{article}

\usepackage[T1]{fontenc}
\usepackage{lmodern}
\usepackage{amsmath,amssymb,amsthm,mathtools}
\usepackage[margin=1in]{geometry}
\usepackage{microtype}
\usepackage[colorlinks=true,linkcolor=blue,urlcolor=blue]{hyperref}

\newtheorem{theorem}{Theorem}
\newtheorem{corollary}[theorem]{Corollary}
\newtheorem{proposition}[theorem]{Proposition}
\theoremstyle{remark}
\newtheorem{remark}[theorem]{Remark}
\newcommand{\R}{\mathbb{R}}
\newcommand{\C}{\mathbb{C}}
\newcommand{\F}{\mathcal{F}}
\newcommand{\supp}{\operatorname{supp}}
\newcommand{\esssupp}{\operatorname*{ess\,supp}}
\newcommand{\one}{\mathbf{1}}

\title{An Isolated Fourier Shell Cannot Support a Nonzero\\
\texorpdfstring{$L^2$}{L2} Bound State}
\author{Richard J. Reyes}
\date{October 26, 2026}

\begin{document}
\maketitle

\begin{abstract}
We prove that a square-integrable function on $\R^d$ whose Fourier
transform is supported on a single wavenumber shell must vanish.  The
argument uses Plancherel's theorem and the zero $d$-dimensional Lebesgue
measure of the shell.  It gives a direct obstruction to nonzero
$L^2$ eigenfunctions of the free-space Helmholtz operator and, more
generally, of nonconstant polynomial functions of the Laplacian.  A
family of modulated wave packets then exhibits the separate freedom
in envelope scale: fixing a carrier wavenumber alone does not select
localization.  The result concerns homogeneous, linear, whole-space
spectral equations; it is not a no-go theorem for nonlinear bound
states or bounded-domain modes.
\end{abstract}

\section{Setting and main theorem}
For $u\in L^2(\R^d;\C)$ we use the unitary Fourier transform
\begin{equation}
 \widehat{u}(\xi)=(2\pi)^{-d/2}
 \int_{\R^d}u(x)e^{-ix\cdot\xi}\,dx,
 \label{eq:Fourier}
\end{equation}
first defined for $u\in L^1\cap L^2$ and then extended by density
and Plancherel to all of $L^2$.  Write
\begin{equation}
 S_{k_0}:=\{\xi\in\R^d:|\xi|=k_0\},\qquad k_0\geq 0.
 \label{eq:shell}
\end{equation}
The support hypothesis below may be read either as a distributional
support condition or as an essential-support condition on the $L^2$
representative; both yield the same conclusion.

\begin{theorem}[Zero $L^2$ state on an isolated Fourier shell]
\label{thm:shell}
Let $d\geq1$, $k_0\geq0$, and $u\in L^2(\R^d;\C)$.  If
\begin{equation}
 \supp\widehat u\subseteq S_{k_0},
 \label{eq:support}
\end{equation}
then
\begin{equation}
 \boxed{u=0\quad\text{almost everywhere on }\R^d.}
 \label{eq:main}
\end{equation}
\end{theorem}

\begin{proof}
By Plancherel's theorem,
\begin{equation}
 \widehat u\in L^2(\R^d),\qquad
 \|u\|_{L^2}^2=\|\widehat u\|_{L^2}^2
 =\int_{\R^d}|\widehat u(\xi)|^2\,d\xi.
 \label{eq:plancherel}
\end{equation}

We next show that $S_{k_0}$ has zero $d$-dimensional Lebesgue measure.
If $d=1$, the shell contains at most the two points $\pm k_0$, so
its Lebesgue measure is zero.  If $d\geq2$ and $k_0>0$, polar
coordinates give
\begin{align}
 |S_{k_0}|
 &=\int_{\R^d}\one_{\{|\xi|=k_0\}}\,d\xi\\
 &=\int_{\mathbb S^{d-1}}\int_0^\infty
     \one_{\{r=k_0\}}r^{d-1}\,dr\,d\omega
 =0.
 \label{eq:null}
\end{align}
For $k_0=0$, the shell is the singleton $\{0\}$, which likewise has
measure zero.  Thus $|S_{k_0}|=0$ in every dimension $d\geq1$.

The support assumption in \eqref{eq:support} implies that
$\widehat u=0$ almost everywhere off $S_{k_0}$.  It follows that
\begin{equation}
 \|\widehat u\|_{L^2}^2
 =\int_{S_{k_0}}|\widehat u(\xi)|^2\,d\xi
 =0,
 \label{eq:zero}
\end{equation}
because $|\widehat u|^2\in L^1(\R^d)$ and the domain of integration
has measure zero.  Equation \eqref{eq:plancherel} now gives
$\|u\|_{L^2}=0$, hence $u=0$ almost everywhere.
\end{proof}

\section{Consequences for homogeneous spectral equations}

\begin{corollary}[Free-space Helmholtz equation]
\label{cor:Helmholtz}
For $d\geq1$ and $k_0\geq0$, the distributional equation
\begin{equation}
 (-\Delta-k_0^2)u=0\quad\text{on }\R^d,
 \qquad u\in L^2(\R^d),
 \label{eq:Helmholtz}
\end{equation}
has only the trivial solution $u=0$.
\end{corollary}

\begin{proof}
Fourier transformation in the sense of tempered distributions yields
\begin{equation}
 (|\xi|^2-k_0^2)\widehat u(\xi)=0
 \quad\text{in }\mathcal S'(\R^d).
 \label{eq:multiplier}
\end{equation}
On the open set $\R^d\setminus S_{k_0}$, the smooth multiplier
$|\xi|^2-k_0^2$ does not vanish.  Local division by that multiplier
in \eqref{eq:multiplier} shows that $\widehat u$ vanishes there as a
distribution.  Consequently
$\supp\widehat u\subseteq S_{k_0}$.  Theorem~\ref{thm:shell}
then gives $u=0$.
\end{proof}

\begin{corollary}[Nonconstant polynomial dispersion]
\label{cor:poly}
Let $P\in\R[t]$ be a nonconstant polynomial, let $\lambda\in\R$,
and let $u\in L^2(\R^d)$ solve
\begin{equation}
 P(-\Delta)u=\lambda u
 \quad\text{in }\mathcal S'(\R^d).
 \label{eq:poly}
\end{equation}
Then $u=0$ almost everywhere.
\end{corollary}

\begin{proof}
The Fourier-side equation is
\begin{equation}
 \bigl(P(|\xi|^2)-\lambda\bigr)\widehat u(\xi)=0.
 \label{eq:poly-mult}
\end{equation}
Since $P(t)-\lambda$ is a nonzero polynomial, it has only finitely
many real roots.  Let $s_1,\ldots,s_N$ denote its nonnegative roots;
the list may be empty.  Local division outside the zero set of the
multiplier shows
\begin{equation}
 \supp\widehat u\subseteq
 \bigcup_{j=1}^N\{\xi:|\xi|=\sqrt{s_j}\}.
 \label{eq:finite-shells}
\end{equation}
Each shell is Lebesgue null by the proof of
Theorem~\ref{thm:shell}.  A finite union of null sets is null; if
$N=0$, the union is empty.  Plancherel therefore implies $u=0$.
\end{proof}

\begin{remark}[General Fourier multipliers]
The same reasoning applies to a Fourier multiplier $m(D)$ at a level
$\lambda$ whenever its eigenvalue equation forces the distributional
support of $\widehat u$ into the set
$\{\xi:m(\xi)=\lambda\}$ and that set has Lebesgue measure zero.
The conclusion need not hold for multipliers that are constant on a
set of positive Lebesgue measure.
\end{remark}

\section{Carrier wavelength versus envelope scale}
The preceding theorem excludes a \emph{nonzero exact-shell $L^2$
eigenstate}.  It does not exclude localized wave packets with nonzero
spectral bandwidth.  The latter also show that selecting a carrier
wavenumber does not independently determine an envelope radius.

\begin{proposition}[Fixed carrier, freely varying envelope]
\label{prop:packet}
Fix $k_0>0$ and a unit vector $\omega_0\in\mathbb S^{d-1}$.
Let $f\in\mathcal S(\R^d)$ be nonzero, with
$\supp\widehat f\subset B(0,\rho)$ for some $\rho>0$.
For each $R>0$ define
\begin{equation}
 u_R(x)=R^{-d/2}f(x/R)e^{ik_0\omega_0\cdot x}.
 \label{eq:packet}
\end{equation}
Then
\begin{align}
 \|u_R\|_2^2&=\|f\|_2^2,\label{eq:packet-mass}\\
 \widehat u_R(\xi)
 &=R^{d/2}\widehat f\bigl(R(\xi-k_0\omega_0)\bigr),
 \label{eq:packet-Fourier}\\
 \supp\widehat u_R
 &\subseteq B(k_0\omega_0,\rho/R),
 \label{eq:packet-support}\\
 \bigl\||x|u_R\bigr\|_2
 &=R\bigl\||x|f\bigr\|_2.
 \label{eq:packet-width}
\end{align}
In particular, a fixed carrier $k_0$ and fixed mass permit a
continuous family of envelope scales $R$.
\end{proposition}

\begin{proof}
For \eqref{eq:packet-mass}, change variables $y=x/R$ and use
$|e^{ik_0\omega_0\cdot x}|=1$.  Equation
\eqref{eq:packet-Fourier} follows from the Fourier scaling and
modulation identities, and \eqref{eq:packet-support} follows at once
from the support of $\widehat f$.  Finally,
\begin{align*}
 \bigl\||x|u_R\bigr\|_2^2
 &=\int_{\R^d}|x|^2R^{-d}|f(x/R)|^2\,dx\\
 &=R^2\int_{\R^d}|y|^2|f(y)|^2\,dy.
\end{align*}
Taking square roots proves \eqref{eq:packet-width}.
\end{proof}

From \eqref{eq:packet-support} and the reverse triangle inequality,
all Fourier components in the packet support satisfy
\begin{equation}
 \bigl||\xi|-k_0\bigr|<\frac{\rho}{R}.
 \label{eq:radial-bandwidth}
\end{equation}
Thus the spectral bandwidth narrows as $R\to\infty$, while the
spatial extent grows linearly in $R$.  A finite-width localized
packet has spectral content beyond an exact shell.  The examples
$u_R$ are kinematic constructions; they are not asserted to solve
any particular evolution or stationary field equation.

\section{Scope and interpretation}
The mathematical conclusion is
\begin{equation}
 \boxed{
 \supp\widehat u\subseteq\{\xi:|\xi|=k_0\},\quad
 u\in L^2(\R^d)
 \quad\Longrightarrow\quad u\equiv0.
 }
 \label{eq:conclusion}
\end{equation}
A nonzero distribution supported on a shell, such as a plane-wave
Fourier delta or a surface measure, can define a generalized
solution, but it is not a nonzero element of $L^2(\R^d)$.

The theorem does \emph{not} rule out nonlinear localized states,
externally trapped states, or modes on bounded and periodic domains.
Those settings do not satisfy the exact whole-space, linear
homogeneous shell-support premise.  Nor does a chosen spectral
minimum, by itself, establish a localized envelope, binding,
stability, or confinement lifetime.  Such conclusions require
additional dynamical or variational hypotheses.

\noindent\textbf{Related research audit:}
\url{https://github.com/rickyjreyes/geometry_of_resonance/blob/main/docs/cross_paper/LOCALIZATION_AND_DOMAIN_SELECTION.md}.

\end{document}
